\documentclass[11pt]{article}
\usepackage[T1]{fontenc}
\usepackage{lmodern,amsmath,amssymb,amsthm,booktabs,microtype}
\usepackage[a4paper,margin=25mm]{geometry}
\usepackage{hyperref}
\hypersetup{
  colorlinks=true,linkcolor=blue,urlcolor=blue,citecolor=blue,
  pdftitle={Minimum-size binary excess coverage arrays},
  pdfsubject={Exact existence and asymptotic width under column symmetry}
}
\newtheorem{theorem}{Theorem}
\newtheorem{lemma}[theorem]{Lemma}
\newcommand{\cx}{\mathsf{CA}_{X}}
\newcommand{\Ksym}{K_{\mathrm{sym}}}
\newcommand{\nint}[1]{\left\lfloor #1\right\rceil}
\newcommand{\code}[1]{\texttt{#1}}
\title{Minimum-size binary excess coverage arrays\\
\large Exact existence and asymptotic width under column symmetry}
\author{AI-assisted research manuscript\\
Author attribution pending human review}
\date{27 November 2026\\
Version 0.2 \enspace\textbullet\enspace Research preprint}
\begin{document}
\maketitle
\begin{abstract}
Gentle, Horsley and Wanless (2025, Section~5) reported that a system
of binomial equations for column-symmetric binary excess coverage arrays
has no nonnegative integer solution at strength \(18\) and width \(36\).
We give a strictly positive solution. All \(37\) multiplicities are
explicit, with minimum \(350\,847\). More generally, every positive
strength \(t\) admits a column-symmetric array with \(2t\) columns and the
minimum possible number of rows, \((t+1)!\). The construction rounds
intersection counts and uses inclusion--exclusion to recover the row
multiplicities. Exact certificates cover the remaining small strengths.
For the largest achievable width under column symmetry we prove the
asymptotic formula \((1+o(1))\log_2((t+1)!)\). Lean~4.19 verifies the
strength-\(18\) certificate and the existence theorem for all positive
strengths. The asymptotic formula is proved mathematically in the text.
\end{abstract}

\section{Problem and main results}
A binary excess coverage array of strength \(t\ge1\) and width \(k\ge t\)
is a multiset of binary words of length \(k\). In each choice of \(t\)
coordinates, each prescribed pattern with \(a\) ones must occur at least
\(a!(t-a)!\) times. Summing these requirements over the patterns in one
choice of coordinates gives the lower bound
\[
  \sum_{a=0}^{t}\binom{t}{a}a!(t-a)! = (t+1)!.
\]
An array with \(N=(t+1)!\) rows therefore realizes each requirement
with equality. This is a binary balanced array with index vector
\(\mu_a=a!(t-a)!\). We call the array \emph{column-symmetric} when
permuting its columns preserves the row multiset. The multiplicity
\(x_i\) is then the number of copies of \emph{each} word of weight \(i\).
Weight-based representations are standard in balanced-array theory
\cite{kuriki,shirakura}.

The constraints for such an array at minimum size are precisely
\begin{equation}
 \sum_{j=0}^{k-t}\binom{k-t}{j}x_{a+j}=a!(t-a)!,
 \qquad 0\le a\le t,\qquad x_i\in\mathbb Z_{\ge0}.                 \label{eq:marg}
\end{equation}
Indeed, the \(k-t\) remaining positions supply the \(j\) additional
ones. Summing the left sides over the \(2^t\) patterns gives the total
number of rows. In the width-\(2t\) case this is the system of Gentle,
Horsley and Wanless \cite[Section~5, equation~(3)]{ghw}. They found
solutions at strengths through \(17\), reported infeasibility for the
symmetric system at \(18\), and left existence without the symmetry
constraint open. The following result settles the instance within
the symmetric class.

\begin{theorem}\label{thm:twice}
For every integer \(t\ge1\), a column-symmetric
\(\cx((t+1)!;t,2t,2)\) exists. The certificate for \(t=18\) has \(37\)
positive coordinates and minimum coordinate \(350\,847\).
\end{theorem}

Let \(\Ksym(t)\) denote the maximum width of a column-symmetric
array with strength \(t\) and exactly \((t+1)!\) rows, for \(t\ge2\).
The following theorem includes finiteness of this maximum.

\begin{theorem}\label{thm:width}
As \(t\to\infty\), \(\Ksym(t)\sim\log_2((t+1)!)\sim t\log_2 t\).
\end{theorem}

\section{Rounding intersection counts}\label{sec:rounding}
Write $\nint{z}$ for a nearest integer to $z$, taking the even integer
at a tie. Define
\begin{equation}
 y_j=\nint{\frac{N}{j+1}},\qquad
 x_i=\sum_{r=0}^{k-i}(-1)^r\binom{k-i}{r}y_{i+r}
 \quad(0\le i,j\le k).                                           \label{eq:mob}
\end{equation}
\begin{lemma}\label{lem:rounding}
The integers in \eqref{eq:mob} satisfy the equalities in
\eqref{eq:marg}. Each $x_i$ is within $2^{k-i}/2$ of
$x_i^*=N/((k+1)\binom{k}{i})$.
\end{lemma}
\begin{proof}
Assign $y_j$ as the total multiplicity of words containing any fixed
$j$-subset in their support. Inclusion--exclusion over the remaining
positions gives \eqref{eq:mob} for the signed multiplicity of a word
with exactly that support. This inversion is valid before positivity
is established. If $j\le t$ then $j+1$ divides $N$, so $y_j=N/(j+1)$.
On a prescribed set of $a$ ones and $b=t-a$ zeros, inclusion--exclusion
therefore gives the signed count
\[
 \sum_{r=0}^{b}(-1)^r\binom{b}{r}\frac{N}{a+r+1}
 = N\int_0^1 p^a(1-p)^b\,dp
 = a!b!.
\]
With every $y_j$ replaced by $N/(j+1)$ the same integral at degree $k$
gives $x_i^*=N i!(k-i)!/(k+1)!$. Each term rounded in \eqref{eq:mob}
contributes error at most half its binomial coefficient.
\end{proof}

\begin{proof}[Proof of Theorem~\ref{thm:twice}]
By Lemma~\ref{lem:rounding}, the condition \(2N\ge(k+1)3^k\) implies
\(x_i\ge0\) for every \(i\), since
\[
  \binom{k}{i}2^{k-i}\le\sum_{j=0}^k\binom{k}{j}2^{k-j}=3^k .
\]
At \(k=2t\), the condition holds at \(t=22\) by integer arithmetic.
Increasing \(t\) by one multiplies the ratio
\(2(t+1)!/((2t+1)3^{2t})\) by
\((t+2)(2t+1)/(9(2t+3))>1\) for \(t\ge22\).
At \(17\le t\le21\), direct evaluation of \eqref{eq:mob} gives the
positive minima recorded in Section~\ref{sec:computation}.
The companion file \code{small.json} supplies the nonnegative integer
vectors at \(1\le t\le16\). The same section describes verification
of the margins and row totals in all these cases. Appendix~\ref{sec:certificate}
lists the full positive solution at \(t=18\).
\end{proof}

The lower bound on the number of rows is attained, so
\(\mathsf{CAN}_X(t,2t,2)=(t+1)!\). In the sequence-covering argument of
\cite[Section~5]{ghw}, absence of this binary array would supply an
obstruction. Theorem~\ref{thm:twice} removes that particular obstruction
at every strength.

\section{Maximum width under column symmetry}\label{sec:width}
Throughout this section \(N=(t+1)!\). We use the convention
\(\binom{n}{r}=0\) when \(n\ge0\) and \(r\notin\{0,\ldots,n\}\).

\begin{lemma}\label{lem:lattice}
For $d=k-t\ge1$, the equalities in \eqref{eq:marg} have an integer
solution $x$ with
\[
 \max_i |x_i-x_i^*| \le \tfrac12\binom{k}{t+1}.
\]
In particular, a nonnegative solution exists whenever
$2N\ge(k+1)\binom{k}{\lfloor k/2\rfloor}\binom{k}{t+1}$.
\end{lemma}
\begin{proof}
There is an integer particular solution $z$. Set $z_0,\ldots,z_{d-1}$
to zero, then solve successively for $z_{d+a}$, whose coefficient
in equation $a$ is one. This also proves that the $t+1$ equations are
independent, so their real kernel has dimension $d$.

For $0\le h<d$, let
\[
 B_{ih}=(-1)^i\binom{i}{h}\binom{k-i}{d-1-h},
 \qquad 0\le i\le k .
\]
Apart from the sign this is a polynomial in $i$ of degree at most
$d-1$. Its alternating $d$th difference vanishes. Thus every column of
$B$ belongs to the kernel. In rows $0,\ldots,d-1$ the entries above
the diagonal vanish and its diagonal entries are nonzero. The columns
are consequently a real basis of the kernel consisting of integer vectors.
The beta-integral calculation from Lemma~\ref{lem:rounding} shows that
\(x^*\) also satisfies the equations, so write $x^*-z=Bv$ and round
the coordinates of $v$. The vector $x=z+B\nint v$ is integral, and the
error in its $i$th coordinate is at most
\[
 \tfrac12\sum_{h=0}^{d-1}\binom{i}{h}\binom{k-i}{d-1-h}
 =\tfrac12\binom{k}{d-1}=\tfrac12\binom{k}{t+1}.
\]
The smallest coordinate of $x^*$ is
$N/((k+1)\binom{k}{\lfloor k/2\rfloor})$, giving the last assertion.
\end{proof}

\begin{proof}[Proof of Theorem~\ref{thm:width}]
For the lower bound take fixed $0<\epsilon<1$ and set
$k=\lfloor(1-\epsilon)\log_2 N\rfloor$. By Stirling's formula,
$k=\Theta(t\log t)$, so $k\ge t+1$ for sufficiently large $t$.
The estimate $\binom{k}{t+1}\le(ek/(t+1))^{t+1}$ gives
$\log((k+1)\binom{k}{t+1})=O(t\log\log t)=o(\log N)$.
Also $\binom{k}{\lfloor k/2\rfloor}\le 2^k\le N^{1-\epsilon}$.
Their product is at most \(N^{1-\epsilon+o(1)}\), which meets
the condition in Lemma~\ref{lem:lattice} for sufficiently large \(t\).

For the upper bound select a uniformly random row from any such array
and write $W$ for its weight. On any $j\le t$ prescribed positions, a
specified pattern of $a$ ones and $b=j-a$ zeros has probability
$\int_0^1p^a(1-p)^b\,dp$. At $j=t$ this is the defining equality.
At smaller $j$ sum over the other $t-j$ coordinates, which multiplies
the integrand by $(p+(1-p))^{t-j}=1$. Conditional on $W=w$, the support
is uniform over the $w$-subsets by column symmetry. For $2s\le t$, the
probability of a specified $s$ ones and $s$ zeros is therefore
\[
 \frac{1}{(2s+1)\binom{2s}{s}}
 =\mathbb E\!\left[\frac{(W)_s(k-W)_s}{(k)_{2s}}\right],
\]
where $(u)_r$ denotes a falling factorial. Some occupied weight $w$
attains at least the displayed average. Using $\binom{2s}{s}\le4^s$
and $(k)_{2s}\ge(k-2s+1)^{2s}$ yields
\[
 \frac{w(k-w)}{k^2}\ \ge\
 \frac{(1-(2s-1)/k)^2}{4(2s+1)^{1/s}}.
\]
With $s=\lfloor\sqrt t\rfloor$ for $t\ge4$, the right side tends to
$1/4$ uniformly in $k\ge t$. Thus $w/k=1/2+o(1)$. The whole layer
of $\binom{k}{w}$ distinct words occurs, so
$\log_2 N\ge\log_2\binom{k}{w}=(1-o(1))k$ by Stirling's formula.
For fixed \(t\ge2\), the earlier probability identity with \(s=1\)
gives an occupied weight \(w\) with \(w(k-w)\ge k(k-1)/6\).
Writing \(m=\min(w,k-w)\) gives \(m\ge(k-1)/6\) and
\(\binom{k}{w}=\binom{k}{m}\ge\binom{2m}{m}\ge2^m\).
Thus \(N\ge2^{(k-1)/6}\), proving finiteness.
Letting \(\epsilon\) decrease to zero finishes the proof.
\end{proof}

\section{Finite verification and reproducibility}\label{sec:computation}
The rounding calculation uses only integers. To form \(y_j\), divide
\(N\) by \(j+1\) with remainder. Round up if twice the remainder exceeds
\(j+1\). At equality, round to the even of the two neighboring
integers. The generator \code{construct.py} then forms consecutive
differences of the list \((y_j)\) to obtain \((x_i)\). The verifier
\code{verify.py} checks the nonnegativity of the supplied multiplicities,
evaluates each sum in \eqref{eq:marg} directly, and independently
checks the row total \(\sum_i\binom{k}{i}x_i\).

\begin{center}
\begin{tabular}{rrrl}
\toprule
Strength \(t\) & Width \(2t\) & \(\min_i x_i\) & Certificate source\\
\midrule
\(1\)--\(16\) & \(2\)--\(32\) & \(\ge 0\) & \code{small.json}\\
17 & 34 & 77\,537 & Rounding construction\\
18 & 36 & 350\,847 & Rounding construction and \code{t18.json}\\
19 & 38 & 1\,792\,479 & Rounding construction\\
20 & 40 & 9\,032\,086 & Rounding construction\\
21 & 42 & 48\,444\,386 & Rounding construction\\
\bottomrule
\end{tabular}
\end{center}

From the \code{excess-arrays/} directory, reproduce the stored and generated
strength-\(18\) certificates with the Python standard library.
\begin{verbatim}
python3 verify.py small.json t18.json
python3 construct.py 18 36 > /tmp/excess-t18.json
python3 verify.py /tmp/excess-t18.json
\end{verbatim}
The strength-\(18\) check reports \(19\) validated equations and row total
\(121\,645\,100\,408\,832\,000\). To regenerate the other rounded cases,
use \code{construct.py} with the strength as its sole argument. The width
defaults to twice the strength. The finite output is in \code{results.csv}
and the build and computation record is in \code{verification.log}.
The repository also includes a SHA-256 manifest.

\section{Machine-checked proof}\label{sec:lean}
The Lean model represents an array by a function from binary words to
natural-number multiplicities. It enumerates the compatible binary
words for each partial pattern and proves that the enumeration has no
duplicates. The total mass of all width-\(k\) words is checked separately.
The existence statement quantifies over each pattern of length \(k\)
with precisely \(k-t\) unspecified positions.

For the small certificates, iterated Pascal sums compute the margins
from the literal vectors of weight-indexed multiplicities. From
strength \(22\) on, the proof rescales the rational counts to integers
and verifies the error estimate and exact margins through identities
for consecutive differences. This gives a proof for arbitrary \(t\),
including strengths outside the finite checks.

\begin{center}
\begin{tabular}{@{}p{.58\linewidth}p{.37\linewidth}@{}}
\toprule
Declaration in namespace \code{Excess} & Mathematical content\\
\midrule
\code{equation\_three\_18} & All \(19\) equations of the original system\\
\code{exists\_binary\_excess\_array\_18} & The concrete width-\(36\) array\\
\code{exists\_symmetric\_binary\_excess\_array} & Existence for every positive strength at width \(2t\)\\
\bottomrule
\end{tabular}
\end{center}

Run \code{lake build} within \code{excess-arrays/lean/} using the pinned
Lean~4.19.0 toolchain. The build prints an axiom audit in
\code{MainResult.lean}. Theorem \code{equation\_three\_18} has an empty
axiom list. The concrete existence theorem uses \code{propext}
and \code{Quot.sound}, and the general existence theorem also uses
\code{Classical.choice}. Theorem~\ref{thm:width} and its asymptotic
estimates are proved in Section~\ref{sec:width}.

\section{Prior work and scope}
Chee, Colbourn, Horsley and Zhou \cite{chee} relate
sequence covering arrays to excess coverage arrays. Gentle, Horsley
and Wanless \cite[Section~5]{ghw} apply that relation to the binary
system considered here. A binary array addresses the resulting
necessary condition. The results here address the binary system in
Section~5. Theorem~1 of \cite{ghw}, which bounds the order of
minimum-size sequence covering arrays of strength \(7\), concerns a
different parameter and construction.

Longyear \cite{longyear}, Shirakura \cite{shirakura} and Kuriki
\cite{kuriki} study general existence criteria and constructions for
balanced arrays. Our claims concern the particular factorial index
vector, the exact row count \(N=(t+1)!\), and, in
Theorem~\ref{thm:width}, column symmetry. The width result gives no
upper bound on arrays outside that symmetry class.

The repository's literature comparison records the journal and
preprint versions, searches for corrections, and searches under both
excess coverage and balanced-array terminology. Those public searches
identified neither a correction for the strength-\(18\) system nor a
prior proof of the symmetric width formula. Priority remains subject
to sources outside the search coverage. This preprint and its
artifacts were prepared with substantial AI assistance in
discovery, proof development, formalization, verification and
writing. Author attribution and independent human review are pending.

\appendix
\section{The strength-18 certificate}\label{sec:certificate}
For each weight \(i\), include \(x_i\) copies of every binary word of
length \(36\) and weight \(i\). The same vector is supplied in
\code{t18.json}.
\begin{center}\small
\begin{tabular}{@{}rr@{\qquad}rr@{\qquad}rr@{}}
\toprule
\(i\) & \(x_i\) & \(i\) & \(x_i\) & \(i\) & \(x_i\)\\
\midrule
0 & 3287704746234543 & 13 & 1711162 & 26 & 12934239\\
1 & 91325561677561 & 14 & 717000 & 27 & 34922343\\
2 & 5218331986327 & 15 & 668613 & 28 & 108647345\\
3 & 460609882625 & 16 & 408459 & 29 & 393846602\\
4 & 55728278943 & 17 & 404361 & 30 & 1687914016\\
5 & 8769561774 & 18 & 350847 & 31 & 8720889081\\
6 & 1660738944 & 19 & 388109 & 32 & 55813690119\\
7 & 408707871 & 20 & 447235 & 33 & 460462943481\\
8 & 100683849 & 21 & 591566 & 34 & 5218580026119\\
9 & 39109461 & 22 & 865672 & 35 & 91325150457081\\
10 & 10767879 & 23 & 1422819 & 36 & 3287705416454919\\
11 & 6580331 & 24 & 2626677 &  & \\
12 & 2062333 & 25 & 5472117 &  & \\
\bottomrule
\end{tabular}
\end{center}

\begin{thebibliography}{9}
\bibitem{ghw}
A.~E. Gentle, D. Horsley, and I.~M. Wanless.
Excess coverage arrays and Levenshtein's conjecture.
\emph{Designs, Codes and Cryptography} \textbf{93} (2025), 5257--5269.
\href{https://doi.org/10.1007/s10623-025-01722-9}{doi:10.1007/s10623-025-01722-9}.
\bibitem{kuriki}
S. Kuriki. On existence and construction of balanced arrays.
\emph{Discrete Mathematics} \textbf{116} (1993), 137--155.
\href{https://doi.org/10.1016/0012-365X(93)90398-D}{doi:10.1016/0012-365X(93)90398-D}.
\bibitem{shirakura}
T. Shirakura. Contributions to balanced fractional $2^m$ factorial designs
derived from balanced arrays of strength $2l$.
\emph{Hiroshima Mathematical Journal} \textbf{7} (1977), 217--285.
\href{https://doi.org/10.32917/hmj/1206135961}{doi:10.32917/hmj/1206135961}.
\bibitem{chee}
Y.~M. Chee, C.~J. Colbourn, D. Horsley, and J. Zhou.
Sequence covering arrays.
\emph{SIAM Journal on Discrete Mathematics} \textbf{27} (2013), 1844--1861.
\href{https://doi.org/10.1137/120894099}{doi:10.1137/120894099}.
\bibitem{longyear}
J.~Q. Longyear. Arrays of strength \(s\) on two symbols.
\emph{Journal of Statistical Planning and Inference} \textbf{10} (1984), 227--239.
\href{https://doi.org/10.1016/0378-3758(84)90075-2}{doi:10.1016/0378-3758(84)90075-2}.
\end{thebibliography}
\end{document}
